% ============================================================
% CHƯƠNG 1: GIỚI THIỆU
% ============================================================
\chapter{GIỚI THIỆU}
\label{Chapter1}

% ------------------------------------------------------------
\section{Bối cảnh nghiên cứu}
\label{sec:boi_canh}
% ------------------------------------------------------------

Chụp X-quang là một kỹ thuật chẩn đoán hình ảnh rất quen thuộc trong khám chữa bệnh vì thao tác nhanh, chi phí không quá cao và có thể quan sát được cấu trúc xương mà không cần can thiệp xâm lấn. Ảnh X-quang thường được dùng để hỗ trợ phát hiện và theo dõi nhiều bất thường như gãy xương, rạn xương, thoái hóa khớp hoặc khối u xương. Lượng ảnh cần xử lý ở các cơ sở y tế cũng không hề nhỏ; chẳng hạn Cambridge University Hospitals ghi nhận hơn 174.000 lượt chụp X-quang mỗi năm, riêng khoa X-quang ngoại trú tiếp nhận khoảng 200--300 người mỗi ngày~\cite{cuh2024xray}.

Trong bài toán phân tích ảnh y khoa, phân đoạn là bước xác định chính xác vùng tổn thương ở cấp độ điểm ảnh. Kết quả phân đoạn có thể hỗ trợ bác sĩ quan sát hình dạng tổn thương, ước lượng kích thước và theo dõi tình trạng theo thời gian. Tuy vậy, nếu làm hoàn toàn thủ công thì công việc này khá tốn thời gian và đòi hỏi sự cẩn thận. Trên ảnh X-quang xương, vấn đề còn khó hơn vì có những tổn thương rất nhỏ, biên không rõ hoặc bị chồng lấp bởi các cấu trúc giải phẫu khác.

Những năm gần đây, học sâu được dùng nhiều trong xử lý ảnh y khoa và cho ra không ít kết quả tốt~\cite{litjens2017survey,shen2017review,minaee2021survey}. Từ các kiến trúc phân đoạn tổng quát như FCN~\cite{long2015fcn}, SegNet~\cite{badrinarayanan2017segnet} và DeepLabv3+~\cite{chen2018deeplabv3plus}, nhiều mô hình hợp với dữ liệu y khoa hơn đã được phát triển như V-Net~\cite{milletari2016vnet}, U-Net~\cite{ronneberger2015unet}, UNet++~\cite{zhou2018unetplusplus} và nnU-Net~\cite{isensee2021nnunet}. Gần đây hơn, Segment Anything Model (SAM)~\cite{2023-SAM-Kirillov} và các biến thể cho ảnh y khoa như SAM-Med2D~\cite{2023-SAMMed2D-Cheng} mở ra hướng phân đoạn dựa trên câu nhắc trực quan, tức là bác sĩ chẩn đoán ảnh có thể cung cấp điểm hoặc hộp giới hạn để định hướng vùng cần xử lý.

Trong khóa luận này, hướng phân đoạn có câu nhắc được chọn vì nó đứng ở khoảng giữa của hai bên: con người vẫn định vị sơ bộ, còn mô hình làm phần khoanh chi tiết ở cấp độ pixel. Cách tiếp cận này thực tế hơn so với việc bắt mô hình tự tìm toàn bộ tổn thương trên ảnh ngay từ đầu, nhất là khi ảnh X-quang xương vốn khó và dễ nhiễu. Nó cũng hợp với hướng làm hệ thống hỗ trợ, tức là mô hình giúp bác sĩ chẩn đoán ảnh làm nhanh hơn và đều tay hơn, chứ không thay hẳn bác sĩ chẩn đoán ảnh~\cite{xu2016deepinteractive,wang2019deepigeos,ma2024medsam}.

% ------------------------------------------------------------
\section{Động lực nghiên cứu}
\label{sec:dong_luc}
% ------------------------------------------------------------

\subsection{Động lực khoa học}

Trong thiết lập phân đoạn tự động, việc bắt mô hình vừa định vị, vừa phân đoạn tổn thương trên ảnh X-quang xương không hề dễ. Nguyên nhân chính là ảnh thường có độ tương phản thấp, nhiều vùng bị chồng cấu trúc và tổn thương có thể rất nhỏ. Khi chuyển sang phân đoạn tương tác, bác sĩ chẩn đoán ảnh có thể cung cấp thêm thông tin không gian, nhưng nếu câu nhắc chỉ được ghép với ảnh ngay từ đầu vào thì về sau tín hiệu này có thể yếu dần qua nhiều tầng tích chập. Vì vậy, một động lực quan trọng của khóa luận là tìm cách đưa câu nhắc vào sâu hơn trong mạng để nó vẫn còn tác dụng ở cả phần mã hóa lẫn giải mã.

Một vấn đề khác cũng rất đáng quan tâm là độ ổn định của mô hình. Trong thực tế, bác sĩ khó vẽ được một hộp giới hạn lúc nào cũng đẹp và bao sát hoàn toàn với tổn thương. Hộp có thể rộng hơn cần thiết hoặc hơi lệch tâm. Nếu mô hình bám quá cứng vào đường biên của hộp thì chỉ cần thao tác lệch một chút là kết quả sẽ giảm rõ rệt. Cho nên khóa luận hướng đến một cách biểu diễn câu nhắc mềm hơn và một cơ chế tiếp nhận linh hoạt hơn để mô hình chịu lỗi tốt hơn trong các tình huống hợp lý.

\subsection{Động lực thực tiễn}

Các mô hình nền tảng dựa trên Transformer có nhiều ưu điểm, nhưng thường nặng nề về số lượng tham số và tài nguyên chạy. Nếu muốn triển khai trên phần cứng hạn chế hoặc trong các tình huống cần phản hồi nhanh thì đây là một trở ngại. So với hướng đó, một mô hình CNN gọn hơn nhưng tận dụng tốt hộp giới hạn đầu vào có lẽ hợp hơn cho một hệ thống hỗ trợ. Trong kịch bản mà khóa luận hướng đến, bác sĩ chỉ cần khoanh sơ bộ vùng nghi ngờ, xem kết quả, rồi nếu cần thì chỉnh lại hộp và chạy lại.

Từ các lý do trên, khóa luận tập trung xây dựng một mô hình phân đoạn ảnh X-quang xương theo ba tiêu chí chính: tận dụng được thông tin từ hộp giới hạn, giữ kết quả tương đối ổn khi câu nhắc bị lệch trong mức hợp lý và có kiến trúc đủ gọn để dễ huấn luyện, cũng như đỡ nặng hơn lúc triển khai.

% ------------------------------------------------------------
\section{Phát biểu bài toán}
\label{sec:phat_bieu}
% ------------------------------------------------------------

Bài toán chính của khóa luận là phân đoạn tương tác vùng tổn thương trên ảnh X-quang xương. Đầu vào của mô hình gồm ảnh X-quang:

\begin{equation}
    \mathbf{I}
    \in
    \mathbb{R}^{H\times W}
\end{equation}

và hộp giới hạn vùng nghi ngờ:

\begin{equation}
    \mathbf{B}
    =
    (x_1,y_1,x_2,y_2),
\end{equation}

trong đó $(x_1,y_1)$ và $(x_2,y_2)$ lần lượt là tọa độ góc trên bên trái và góc dưới bên phải. Hộp giới hạn được chuyển thành bản đồ nhiệt câu nhắc:

\begin{equation}
    \mathbf{H}(\mathbf{B})
    \in
    [0,1]^{H\times W}.
\end{equation}

Mục tiêu của mô hình là dự đoán mặt nạ nhị phân:

\begin{equation}
    \widehat{\mathbf{M}}
    \in
    \{0,1\}^{H\times W},
\end{equation}

trong đó giá trị 1 biểu thị điểm ảnh được dự đoán thuộc vùng tổn thương. Quá trình suy luận được biểu diễn bởi:

\begin{equation}
    \widehat{\mathbf{M}}
    =
    \mathbf{1}
    \left[
        f_{\mathrm{PGA}}
        \left(
            \mathbf{I},
            \mathbf{H}(\mathbf{B});
            \boldsymbol{\theta}
        \right)
        \geq 0{,}5
    \right],
    \label{eq:problem_segmentation}
\end{equation}

trong đó $f_{\mathrm{PGA}}$ là PGA-UNet,
$\boldsymbol{\theta}$ là bộ tham số học được và
$\mathbf{1}[\cdot]$ là phép nhị phân hóa đầu ra xác suất tại ngưỡng 0,5.

Mô hình được kỳ vọng đáp ứng các yêu cầu sau:

\begin{itemize}
    \item Dự đoán mặt nạ có mức độ chồng lấp cao với nhãn tham chiếu;
    \item Duy trì đường biên và vị trí vùng tổn thương;
    \item Duy trì hiệu năng tốt khi hộp giới hạn rộng hơn hoặc bị lệch khỏi vị trí lý tưởng;
    \item Có số lượng tham số và chi phí suy luận phù hợp với quá trình tương tác.
\end{itemize}

Trong giai đoạn huấn luyện, hộp giới hạn được tạo từ mặt nạ nhãn gốc rồi biến đổi theo một số kịch bản mô phỏng. Khi suy luận, hộp sẽ do bác sĩ cung cấp. Do đó, kết quả thực nghiệm trong khóa luận phản ánh hành vi của mô hình trong thiết lập đã có thông tin định vị sơ bộ khá gần với tổn thương, chứ chưa thể xem là đánh giá đầy đủ cho trường hợp bác sĩ vẽ tự do trong thực tế. Phần đánh giá với câu nhắc thực tế được xem là hướng phát triển tiếp theo.

Ngoài bài toán phân đoạn chính, khóa luận còn có một thử nghiệm mở rộng với Gatekeeper dùng EfficientNet-B3. Thành phần này chỉ đóng vai trò gợi ý sàng lọc ảnh trước khi bác sĩ xác nhận và cung cấp hộp giới hạn. Nó không phải là đóng góp kiến trúc chính của khóa luận và cũng không thay thế quyết định của bác sĩ chẩn đoán ảnh.

% ------------------------------------------------------------
\section{Thách thức của bài toán}
\label{sec:thach_thuc}
% ------------------------------------------------------------

Việc xây dựng một mô hình phân đoạn tương tác cho ảnh X-quang xương gặp bốn nhóm khó khăn chính:

\begin{itemize}
    \item \textbf{Dữ liệu và chú thích:}
    Mặt nạ phân đoạn ở cấp độ điểm ảnh tốn nhiều thời gian để xây dựng, trong khi số lượng ảnh bệnh lý có chú thích thường hạn chế. Ngoài ra, một ảnh có thể chứa nhiều vùng tổn thương được biểu diễn bởi các đa giác riêng biệt, đòi hỏi quy trình tiền xử lý và đánh giá nhất quán.
    \item \textbf{Đặc điểm ảnh X-quang:}
    Ảnh X-quang có độ tương phản không đồng đều và chứa nhiều cấu trúc giải phẫu chồng lấp. Vùng khối u có thể có hình dạng bất thường, trong khi đường gãy thường nhỏ, mảnh và khó phân biệt với các đường biên xương tự nhiên.

    \item \textbf{Tích hợp và duy trì câu nhắc:}
    Mô hình cần khai thác thông tin không gian từ hộp giới hạn tại nhiều mức đặc trưng. Câu nhắc không nên chỉ ảnh hưởng tại đầu vào rồi suy giảm qua các tầng sâu, đồng thời cũng không được khiến mô hình phụ thuộc cứng vào đường biên hộp.

    \item \textbf{Độ ổn định và hiệu quả tính toán:}
    Mô hình cần duy trì hiệu năng khi câu nhắc bị lệch tâm hoặc thay đổi kích thước. Bên cạnh độ chính xác, kiến trúc phải đủ gọn để hỗ trợ suy luận lặp lại khi bác sĩ điều chỉnh hộp giới hạn.
\end{itemize}

% ------------------------------------------------------------
\section{Đóng góp của khóa luận}
\label{sec:dong_gop}
% ------------------------------------------------------------

Đóng góp chính của khóa luận là đề xuất và đánh giá PGA-UNet, một kiến trúc CNN hạng nhẹ cho bài toán phân đoạn tương tác ảnh X-quang xương. Cụ thể hơn, khóa luận có các đóng góp sau:

\begin{enumerate}
    \item \textbf{Biểu diễn hộp giới hạn bằng bản đồ nhiệt dạng cao nguyên.}

    Hộp giới hạn được chuyển thành một bản đồ hai chiều có giá trị cao tương đối đều ở bên trong và được làm mềm ở phần biên bằng bộ lọc Gaussian. Cách biểu diễn này giúp mô hình không chỉ chú ý vào tâm hộp mà còn giữ được thông tin trên toàn bộ vùng câu nhắc. Hiệu quả của dạng Gaussian và dạng nhị phân được so sánh lại trong phần nghiên cứu loại bỏ thành phần.
    \item \textbf{Tích hợp câu nhắc vào bộ mã hóa bằng PSG.}

    Cổng không gian PSG (\textit{Prompt Spatial Gate}) dùng bản đồ nhiệt để điều biến đặc trưng ở các tầng mã hóa. Có thể hiểu đơn giản là mô hình được nhắc phải chú ý nhiều hơn vào vùng mà bác sĩ đã khoanh, nhưng vẫn giữ đường truyền của đặc trưng ảnh gốc nhờ kết nối dư.

    \item \textbf{Điều kiện hóa bộ giải mã bằng CAD.}

    CAD (\textit{Conditional Attention Decoder}) mở rộng cổng chú ý của Attention U-Net~\cite{oktay2018attentionunet} bằng cách đưa thêm đặc trưng câu nhắc vào tín hiệu điều khiển. Nhờ đó, phần chú ý ở kết nối tắt không chỉ phụ thuộc vào đặc trưng ảnh mà còn phụ thuộc vào vùng bác sĩ đang quan tâm.

    \item \textbf{Đánh giá thiết kế trên hai bộ dữ liệu và nhiều điều kiện câu nhắc.}

    PGA-UNet được huấn luyện và đánh giá riêng trên BTXRD và FracAtlas, là hai bộ dữ liệu có dạng tổn thương khác nhau gồm khối u xương và gãy xương. Thực nghiệm bao gồm so sánh với U-Net, Attention U-Net và SAM-Med2D; đánh giá trong các kịch bản câu nhắc Bao trọn, Lệch tâm và Hỗn hợp; nghiên cứu loại bỏ thành phần; kiểm chứng chéo và phân tích theo từng nhóm tổn thương.

    \item \textbf{Khảo sát khả năng tích hợp vào luồng xử lý hỗ trợ.}

    Khóa luận bổ sung Gatekeeper như một thử nghiệm mở rộng để minh họa cách PGA-UNet có thể đặt trong một luồng xử lý gồm cả ảnh bình thường và ảnh bệnh lý. Thành phần này chỉ có ý nghĩa hỗ trợ sàng lọc, không được xem là đóng góp kiến trúc độc lập.
\end{enumerate}

Trong các nội dung trên, phần cốt lõi nhất vẫn là cách phối hợp giữa biểu diễn câu nhắc, PSG và CAD. Hiệu quả của từng phần sẽ được kiểm tra lại bằng thực nghiệm ở Chương~\ref{Chapter4}.
% ------------------------------------------------------------
\section{Phạm vi nghiên cứu}
\label{sec:pham_vi}
% ------------------------------------------------------------

Khóa luận giới hạn trong phạm vi sau:

\begin{itemize}
    \item Nghiên cứu phân đoạn ảnh X-quang xương hai chiều;
    \item Sử dụng ảnh bệnh lý có mặt nạ cấp độ điểm ảnh để huấn luyện mô-đun phân đoạn;
    \item Sử dụng hộp giới hạn làm dạng câu nhắc trực quan chính;
    \item Đánh giá trên hai bộ dữ liệu công khai BTXRD và FracAtlas;
    \item Mô phỏng sai lệch câu nhắc bằng các phép mở rộng và dịch chuyển hộp giới hạn theo quy tắc có kiểm soát, không thay thế cho đánh giá với câu nhắc thực do bác sĩ cung cấp;
    \item Chưa thực hiện thử nghiệm với câu nhắc do bác sĩ vẽ trực tiếp hoặc dữ liệu thô được thu thập tại cơ sở y tế;
    \item Gatekeeper chỉ được khảo sát như một thành phần hỗ trợ, không phải hệ thống chẩn đoán tự động.
\end{itemize}

Các kết quả trong khóa luận chỉ phản ánh hiệu năng kỹ thuật trong điều kiện thực nghiệm có kiểm soát. Chúng không được dùng như kết luận chẩn đoán và cũng chưa đủ để khẳng định hệ thống đã sẵn sàng cho triển khai lâm sàng.
% ------------------------------------------------------------
\section{Bố cục khóa luận}
\label{sec:bo_cuc}
% ------------------------------------------------------------

Khóa luận được trình bày trong năm chương:

\begin{itemize}
    \item \textbf{Chương 1 -- Giới thiệu:}
    Trình bày bối cảnh, động lực, bài toán nghiên cứu, các thách thức, đóng góp và phạm vi của khóa luận.

    \item \textbf{Chương 2 -- Các công trình liên quan và cơ sở lý thuyết:}
    Khảo sát phân đoạn ảnh y khoa tự động, phân đoạn tương tác và mô hình nền tảng dựa trên câu nhắc; trình bày cơ sở về U-Net, cơ chế chú ý, biểu diễn câu nhắc, điều biến đặc trưng và các độ đo đánh giá.

    \item \textbf{Chương 3 -- Mô hình và hệ thống đề xuất:}
    Trình bày kiến trúc PGA-UNet, gồm bản đồ nhiệt câu nhắc dạng cao nguyên, PSG và CAD; đồng thời mô tả quy trình huấn luyện, suy luận và luồng xử lý hỗ trợ với Gatekeeper.

    \item \textbf{Chương 4 -- Thực nghiệm và đánh giá:}
    Trình bày dữ liệu, cấu hình huấn luyện, kết quả so sánh với các mô hình cơ sở và SAM-Med2D, mức độ ổn định dưới sai lệch câu nhắc mô phỏng, tóm tắt các kết luận chính về đóng góp kiến trúc và tính nhất quán trên hai bộ dữ liệu, cùng tóm tắt thử nghiệm mở rộng với Gatekeeper. Toàn bộ bảng số liệu chi tiết của nghiên cứu loại bỏ thành phần, kiểm chứng chéo và phân tích theo đặc tính tổn thương được trình bày tại Phụ lục~\ref{AppendixA}; số liệu chi tiết của luồng xử lý hai giai đoạn với Gatekeeper được trình bày tại Phụ lục~\ref{AppendixB}.

    \item \textbf{Chương 5 -- Kết luận và hướng phát triển:}
    Tổng hợp các kết quả đạt được, chỉ ra hạn chế của nghiên cứu và đề xuất các hướng phát triển tiếp theo.
\end{itemize}
